% ============================================================
%  chapters/03_literature_review.tex
% ============================================================

\section{Literature Review}
\label{sec:literature}

Given the role education plays in poverty reduction, productivity growth,
and development, questions about the effectiveness of educational investments
are of fundamental importance to policymakers in both developed and
developing countries. One issue that has received sustained attention in this
literature is school dropout---or early exit from school---and its
determinants.

Several factors beyond simple inability to pay have been shown to drive
dropout. Lack of school quality, or the \emph{perceived} lack thereof,
plays an important role \citep{levy1971, hanushek2008, patron2008}.
\citet{hanushek2008}, studying primary school-age children in Egypt, find
that students are more likely to drop out when attending low-quality schools,
holding all other factors constant. \citet{patron2008} similarly finds that
perceived quality is an important determinant of secondary school attendance
in Uruguay. \citet{levy1971} points to additional factors including high
grade repetition rates, elevated social tension, and high fertility rates
as negative drivers of educational attainment. For older primary students,
the economic return to schooling and its opportunity cost also matter: as
\citet{ohiggins2008} show in the case of Italy, dropout decisions can be
entirely rational and may occur even in the absence of binding monetary
constraints.

The primary concern of this paper, however, is the relationship between
violence and educational attainment---a topic that has been widely studied
and which generally points in the same direction. Exposure to violence has
been found to have a negative effect on educational attainment across a
range of settings and conflict types \citep{akresh2008, swee2009,
rodriguez2009, shemyakina2011, leon2012}.

\citet{akresh2008} study the Rwandan genocide and compare the educational
outcomes of children exposed to the violence with those who were not.
They find that exposed children achieved 0.5 fewer years of schooling and
were 15 percent less likely to complete third grade. \citet{shemyakina2011}
examines the civil war in Tajikistan and finds that school-age girls
exposed to violence were the most negatively affected cohort relative to
unexposed girls and boys of both exposure statuses. \citet{leon2012}
finds that children exposed to violence at a very young age---before
starting school---attained 0.31 fewer years of education than unexposed
children, and that early exposure has a persistent effect, whereas older
children exposed to violence after starting school eventually catch up
with their non-exposed peers.

Not all studies find uniformly negative effects at all margins. \citet{swee2009},
studying the Bosnian war, finds that greater war intensity reduced the
probability of completing secondary school but had no effect on primary
school completion rates. He attributes reduced secondary enrollment to
demand-side effects, specifically the increased likelihood of military
conscription. \citet{rodriguez2009} examine Colombian children aged 6--17
and find that children older than 11 are most affected by conflict through
induced early labor market entry, mediated by increased mortality risk,
negative income shocks, and reduced school quality.

Taken together, this literature establishes that the relationship between
violence and schooling is negative in expectation but heterogeneous in
magnitude across age of exposure, gender, conflict type, and outcome
measured. This paper contributes to the literature by applying these
frameworks to the Nigerian context, where violence has been pervasive but
where large-scale, nationally representative education data offer an
opportunity to estimate effects across a diverse population spanning multiple
ethnic, religious, and geographic contexts.
